\section{Future Work} \label{sec:FUTURE_WORK}
% have lots of problem vs. solutions pairs

\todo{andere tools anpassen, z.b. explorer, inspector, style editor,
  text editor, (parts bin - group aware)}

A number of open questions remain in our approach, that need to be
explored in future work.

\paragraph{Groups as Second-Class Objects}
There is no object representation of a group, instead, group
membership is declared by tagging member objects. The list of groups
for a world is derived from the set of tags on the live
objects. Selecting a group is akin to selecting all live objects that
have the group tag. Tags are persisted in the PartsBin, so an object
that is saved, removed from the world, and later restored is still a
member of the group. However, if the group was edited between removing
and restoring a member of that group, that member will not receive the
changes, and the group members will be inconsistent with each other.

\paragraph{Similarity Based Groups}
Selection of a group presents the intersection of all members'
properties and functions. This presents two problems: first, by adding
properties to a group developers may accidentally overwrite properties
that some, but not all of the members have. Second, there may be
sub-groups that share more behavior than the whole group, but finding
those currently involves checking all groups in $\mathcal{P}(M)$ for a
set of objects $M$ manually.

A possible solution to these two problems may be to automatically
generate anonymous sub-groups that encapsulate the differences between
group members. Sub-groups allow developers to identify diverging group
members and eventually merge them, as well recognize natural
sub-groups within groups, which they can then name appropriately.

\paragraph{Method Ownership for Groups}
Given an object that is part of a number of groups, our approach
offers no way of knowing whether a function is part of a group, or
unique to this object. This means developers cannot know whether they
will evict this object from a group by editing a function directly. A
possible solution may be to tag functions with group ownership as well
as objects. If a function was edited as part of a group edit, the
group tag for that function is changed to point to the owning
group. If developers decide to edit a function on a single object that
was part of a group, the annotation is removed.

%%% Local Variables: 
%%% mode: latex
%%% TeX-master: "../paper"
%%% End: 
